\chapter{Reproduction du papier}

\section{Objet et perimetre de la reproduction}

L'objectif de cette partie est de reproduire, aussi fidelement que possible,
les resultats presentes dans \emph{Dial-In LLM: Human-Aligned LLM-in-the-loop
Intent Clustering for Customer Service Dialogues}~\cite{hong2025dialinllm}.
Compte tenu des ressources publiques disponibles, la reproduction a ete
restreinte a la partie anglaise du protocole, c'est-a-dire aux benchmarks
Bank77, CLINC(I), MTOP(I) et Massive(I).

Ce choix est methodologiquement justifie par trois elements. En premier lieu,
les jeux de donnees concernes sont publics et leurs tailles sont explicites
dans le papier. En second lieu, les scores cibles sont rapportes de maniere
precise par les auteurs, ce qui permet une comparaison quantitative directe.
En troisieme lieu, cette partie du protocole mobilise deja l'ensemble des
composantes essentielles du pipeline \emph{LLM-in-the-loop}. En revanche, la
reproduction stricte des resultats principaux sur le corpus chinois n'est pas
realiste dans notre environnement, en raison de l'absence des checkpoints
fine-tunes utilises par les auteurs.

Les scores de reference extraits du papier sont rappeles dans le
tableau~\ref{tab:paper-targets}.

\begin{table}[h]
  \centering
  \begin{tabular}{lr}
    \toprule
    Benchmark & NMI papier (\%) \\
    \midrule
    Bank77 & 82.32 \\
    CLINC(I) & 94.12 \\
    MTOP(I) & 72.45 \\
    Massive(I) & 78.12 \\
    \bottomrule
  \end{tabular}
  \caption{Scores de reference rapportes dans le papier pour les benchmarks anglais.}
  \label{tab:paper-targets}
\end{table}

\section{Description de l'approche etudiee}

L'approche proposee dans \emph{Dial-In LLM} repose sur une boucle iterative de
clustering dans laquelle un grand modele de langage intervient comme composant
de decision. Le pipeline peut etre decompose en cinq etapes principales :
\begin{enumerate}[leftmargin=*]
  \item encodage des enonces dans un espace vectoriel ;
  \item generation de plusieurs partitions candidates a partir de valeurs
  concurrentes de \texttt{K} ;
  \item echantillonnage representatif des clusters obtenus ;
  \item evaluation de la coherence des clusters par le LLM, avec generation
  d'un jugement \texttt{Good/Bad} et d'un nom d'intention ;
  \item fusion finale de clusters semantiquement proches.
\end{enumerate}

Dans notre implementation, ce schema general a ete conserve. La boucle
iterative, la recherche de \texttt{K}, l'evaluation de coherence, le nommage
des clusters et le merge final ont ete reconstruits au sein d'un pipeline
entierement executable. En revanche, plusieurs composantes ont du etre
approximees :
\begin{itemize}[leftmargin=*]
  \item le LLM de reference du papier pour les benchmarks anglais n'etait pas
  publiquement accessible ;
  \item le \texttt{convex sampling} a d'abord ete remplace par des strategies
  plus simples, puis par une approximation locale ;
  \item la procedure de fusion finale reste plus simple que la version
  probabiliste decrite dans le papier.
\end{itemize}

La reproduction obtenue doit donc etre comprise comme une reproduction
approchee, documentee et comparee explicitement au protocole original.

\section{Protocole de reproduction}

\subsection{Alignement des donnees}

Une phase preliminaire d'alignement a ete necessaire afin de garantir que les
experiences portent sur les memes objets d'evaluation que dans le papier. Les
benchmarks utilises et leurs tailles observees sont resumes dans le
tableau~\ref{tab:benchmark-sizes}.

\begin{table}[h]
  \centering
  \begin{tabular}{lrr}
    \toprule
    Benchmark & Taille observee & Taille papier \\
    \midrule
    Bank77 & 3080 & 3080 \\
    CLINC(I) & 4500 & 4500 \\
    MTOP(I) & 4386 & 4386 \\
    Massive(I) & 2974 & 2974 \\
    \bottomrule
  \end{tabular}
  \caption{Verification des tailles des benchmarks utilises pour la reproduction.}
  \label{tab:benchmark-sizes}
\end{table}

Cet alignement a implique :
\begin{itemize}[leftmargin=*]
  \item l'export explicite des splits \texttt{test} dans le format attendu par
  le pipeline ;
  \item le filtrage \texttt{intent-only} pour CLINC(I) ;
  \item la desactivation de la deduplication interne afin d'eviter toute
  modification artificielle des benchmarks ;
  \item la sauvegarde systematique de la configuration de run pour garantir la
  tracabilite des experiences.
\end{itemize}

\subsection{Configuration retenue}

Le papier indique que les benchmarks anglais sont traites avec un
\texttt{Llama-7B} fine-tune sur 800 clusters d'intentions annotes. Cette
configuration n'etait pas directement accessible dans notre environnement. La
configuration suivante a donc ete retenue comme approximation la plus proche
effectivement executable :
\begin{itemize}[leftmargin=*]
  \item embeddings : \texttt{BAAI/bge-large-en-v1.5} ;
  \item evaluateur et generateur de noms :
  \texttt{mistralai/Mistral-7B-Instruct-v0.3} ;
  \item algorithme de clustering : \texttt{MiniBatchKMeans} ;
  \item \texttt{sample-size = 10} ;
  \item \texttt{tmax = 5} ;
  \item recherche par candidats de \texttt{K} ;
  \item \texttt{dedupe = false}.
\end{itemize}

Les principaux ecarts assumes par rapport au protocole du papier sont donc :
\begin{itemize}[leftmargin=*]
  \item l'absence du \texttt{Llama-7B} fine-tune utilise par les auteurs ;
  \item l'absence du \texttt{convex sampling} exact ;
  \item une simplification partielle de la logique de merge ;
  \item un prompting adapte au LLM local retenu.
\end{itemize}

\subsection{Justification du non-recours a \texttt{Llama-7B} et au \texttt{convex sampling} exact}

Le non-recours au \texttt{Llama-7B} mentionne dans le papier ne releve pas d'un
choix arbitraire, mais d'une contrainte de reproductibilite. Les auteurs
indiquent utiliser, pour les benchmarks anglais, un \texttt{Llama-7B}
fine-tune sur 800 clusters d'intentions annotes. Or ce checkpoint fine-tune
n'etait pas publiquement disponible dans le cadre de ce projet. En
consequence, il n'etait pas possible de reproduire strictement cette
composante du protocole. Le recours a
\texttt{mistralai/Mistral-7B-Instruct-v0.3} doit donc etre compris comme une
approximation operationnelle, choisie parce qu'elle etait accessible,
compatible avec l'environnement serveur et suffisamment stable pour soutenir
une campagne de reproduction complete.

La meme logique s'applique au \texttt{convex sampling}. Le papier en souligne
l'importance, mais ni une implementation directement reutilisable ni une
description suffisante pour garantir une reconstruction exacte n'etaient
disponibles dans les ressources publiques exploitees ici. Dans un premier
temps, des strategies plus simples et plus robustes ont donc ete retenues pour
stabiliser le pipeline. Dans un second temps, une approximation locale du
\texttt{convex sampling} a ete introduite a titre d'ablation. Les resultats
obtenus n'ont toutefois pas montre d'amelioration mesurable, ce qui a conduit
a ne pas retenir cette variante dans la configuration de reference.

\section{Environnement experimental}

L'ensemble des experiences a ete execute sur le serveur Jupyter du projet, dans
le repertoire \texttt{/nfs/LLM-DIAL}. Le cadre d'execution a ete fixe de facon
stricte afin de garantir la reproductibilite des runs et de limiter les effets
de bord.

\subsection{Plateforme d'execution}

Les calculs ont ete menes sur une machine disposant de deux GPU NVIDIA L40. Par
convention, toutes les experiences ont ete executees sur le second GPU
physique, afin de conserver un protocole d'execution uniforme sur l'ensemble du
projet. Les caches de modeles ont ete confines au depot local du projet, de
meme que l'ensemble des artefacts de sortie.

\subsection{Environnement logiciel}

Au moment des principales experiences de reproduction, l'environnement
logiciel etait le suivant :
\begin{itemize}[leftmargin=*]
  \item Python 3.12.6 ;
  \item \texttt{torch} disponible ;
  \item \texttt{transformers} disponible ;
  \item \texttt{sentence-transformers} disponible ;
  \item \texttt{accelerate} indisponible ;
  \item \texttt{bitsandbytes} indisponible.
\end{itemize}

Ces caracteristiques ont directement conditionne le niveau de fidelite
accessible. En particulier, l'absence d'outils de quantification et
l'impossibilite de charger le modele exact des auteurs ont impose le recours a
une configuration alternative documentee.

\section{Resultats de reproduction}

La campagne de reproduction principale a produit les resultats presentes dans
le tableau~\ref{tab:reproduction-results}.

\begin{table}[h]
  \centering
  \small
  \begin{tabular}{lrrrrrr}
    \toprule
    Benchmark & NMI papier & NMI reproduit & Delta & Couverture & Clusters & Restant \\
    \midrule
    Bank77 & 82.32 & 83.78 & +1.46 & 98.77 & 142 & 38 \\
    CLINC(I) & 94.12 & 90.90 & -3.22 & 95.96 & 233 & 182 \\
    MTOP(I) & 72.45 & 71.12 & -1.33 & 96.26 & 185 & 164 \\
    Massive(I) & 78.12 & 72.44 & -5.68 & 87.56 & 104 & 370 \\
    \bottomrule
  \end{tabular}
  \caption{Comparaison entre les scores du papier et la baseline de reproduction.}
  \label{tab:reproduction-results}
\end{table}

Ces resultats montrent que la reproduction obtenue est deja competitive sur
plusieurs benchmarks. Le score atteint sur Bank77 depasse legerement celui
rapporte par les auteurs, tandis que MTOP(I) reste proche de la cible. En
revanche, les ecarts demeurent plus importants sur CLINC(I) et Massive(I), ce
qui signale des limites plus structurelles du pipeline reproduit ou de son
approximation locale.

Un second enseignement concerne la structure meme des partitions produites. La
couverture reste elevee sur trois benchmarks, mais le nombre de clusters finaux
demeure souvent superieur au nombre de classes de reference. Cette observation
oriente l'analyse vers une tendance a la sur-fragmentation, qui deviendra un
axe central des ameliorations ulterieures.

\section{Analyse critique de l'approche et des evaluations proposees par les auteurs}

\subsection{Points forts de l'approche}

La reproduction permet de confirmer plusieurs qualites du papier. La premiere
est l'integration effective du LLM dans la boucle de clustering. Cette
integration ne se reduit pas a un simple enrichissement des representations :
le modele intervient dans la selection des clusters retenus et dans la
construction d'une couche d'interpretabilite via le nommage des intents.

La deuxieme qualite tient au caractere iteratif de la methode. Le nombre final
de clusters n'est pas suppose connu a l'avance, ce qui correspond mieux aux
conditions reelles d'exploration de corpus conversationnels.

La troisieme qualite est d'ordre experimental : le papier propose une
evaluation sur plusieurs benchmarks anglais, ainsi qu'un corpus chinois de
grande ampleur, ce qui donne un cadre comparatif plus riche que celui de
travaux limites a un seul jeu de donnees.

\subsection{Limites de l'approche}

La reproduction met egalement en evidence plusieurs limites.

La premiere concerne la dependance au LLM utilise dans la boucle. Le
comportement du pipeline est sensible a la qualite des jugements
\texttt{Good/Bad} et au nommage des clusters ; toute difference de modele peut
donc affecter significativement les resultats.

La deuxieme tient au role du sampling, du merge et de la recherche de
\texttt{K}. Ces composantes ont un impact structurel sur le niveau de
fragmentation final, mais le papier ne permet pas toujours de les reconstruire
integralement a partir des seules ressources publiques.

La troisieme limite est d'ordre computationnel. Le recours repete au LLM dans
la boucle de clustering augmente le cout d'execution et complique la
reproductibilite stricte lorsque les memes ressources logicielles ou materielles
ne sont pas disponibles.

\subsection{Discussion du protocole d'evaluation des auteurs}

Le protocole d'evaluation des auteurs presente un interet indeniable. L'usage
combine de la NMI et d'un \emph{goodness score} permet de distinguer la qualite
externe de la partition et sa coherence semantique interne. Ce choix est
methodologiquement pertinent dans le cadre du clustering d'intentions.

Toutefois, plusieurs limites doivent etre soulignees. D'abord, la NMI, bien
qu'utile, ne capture pas a elle seule l'ensemble de la qualite semantique des
clusters. Ensuite, l'absence des checkpoints fine-tunes utilises dans le papier
rend difficile une reproduction strictement equivalente de l'evaluation
rapportee. Enfin, les resultats obtenus dans notre reproduction montrent que la
reaction a une meme configuration varie sensiblement selon les benchmarks, ce
qui plaide pour une lecture benchmark par benchmark plutot que pour une
interpretation globalisante.

\section{Bilan}

La reproduction obtenue constitue une base experimentale solide. Elle permet de
valider le coeur de la methode dans un cadre realisticement reproductible, tout
en localisant de maniere argumentee les principales sources d'ecart avec le
papier. Cette baseline servira de point de comparaison pour l'ensemble des
evaluations complementaires et des modifications methodologiques presentees
dans le chapitre suivant.
